%------------------------------------------------------------------------------%
\documentclass[a4paper,12pt,fleqn]{article}

\usepackage{gaal}

%\ShowAnswers

%------------------------------------------------------------------------------%
\begin{document}

\makeHeader{GAAL}{Prova 1 -- Turma 4}{26/08/2026}

% 01
%------------------------------------------------------------------------------%
\exercise{25\%}
Dados
\[
  A = \begin{bmatrix}
    1 & 2 & -1 \\
    0 & 3 & 2
  \end{bmatrix}
\qquad
  B = \begin{bmatrix}
    2 & -1 & 0 \\
    1 & 1  & 3
  \end{bmatrix}
\qquad
  C = \begin{bmatrix}
    1 & 2  & 3 \\
    0 & -1 & 1
  \end{bmatrix}
\]
calcule \(AC^t + BC^t\)

\begin{answer}
  \[
    AC + BC
    = (A + B)C
  \]
  \begin{align*}
    A + B
    & =
      \begin{bmatrix}
        1 & 2 & -1 \\
        0 & 3 & 2
      \end{bmatrix}
      +
      \begin{bmatrix}
        2 & -1 & 0 \\
        1 & 1  & 3
      \end{bmatrix}
    \\[2mm]
    & =
      \begin{bmatrix}
        1 + 2 & 2 - 1 & -1 + 0 \\
        0 + 1 & 3 + 1 &  2 + 3
      \end{bmatrix}
    \\[2mm]
    & =
      \begin{bmatrix}
        3 & 1 & -1 \\
        1 & 4 &  5
      \end{bmatrix}
  \end{align*}
  \begin{align*}
    AC + BC
    & = (A + B)C
    \\
    & =
      \begin{bmatrix}
        3 & 1 & -1 \\
        1 & 4 &  5
      \end{bmatrix}
      \begin{bmatrix}
        1 & 0  \\
        2 & -1 \\
        3 & 1
      \end{bmatrix}
    \\[2mm]
    & =
      \begin{bmatrix}
        3 \cdot 1 + 1 \cdot 2 + (-1) 3 &
        3 \cdot 0 + 1 (-1)    + (-1) 1
        \\
        1 \cdot 1 + 4 \cdot 2 +  5 \cdot 3 &
        1 \cdot 0 + 4 (-1)    +  5 \cdot 1
      \end{bmatrix}
    \\[2mm]
    & =
      \begin{bmatrix}
        3 + 2 -  3 & 0 - 1 - 1 \\
        1 + 8 + 15 & 0 - 4 + 5
      \end{bmatrix}
    \\[2mm]
    & =
      \begin{bmatrix}
        2 & -2 \\
        24 & 1
      \end{bmatrix}
  \end{align*}
\end{answer}

% 02
%------------------------------------------------------------------------------%
\exercise{25\%}
Resolva o sistema linear por Gauss--Jordan
\[
  \systeme[sort={x, y, z, w}]{
     x + 2y -  z + w = 3,
    2x + 5y +      w = 7,
    -x -  y + 2z - w = -1
  }
\]

\begin{answer}
Matriz aumentada
\[
\begin{amatrix}{4}
  1  & 2  & -1 & 1  & 3  \\
  2  & 5  & 0  & 1  & 7  \\
  -1 & -1 & 2  & -1 & -1
\end{amatrix}
\]
Utilizamos a primeira linha para eliminar os elementos abaixo
do primeiro pivô
\[
  L_2\leftarrow L_2-2L_1
\]
\[
  L_3\leftarrow L_3+L_1
\]
Assim
\[
\begin{amatrix}{4}
  1 & 2 & -1 & 1  & 3 \\
  0 & 1 & 2  & -1 & 1 \\
  0 & 1 & 1  & 0  & 2
\end{amatrix}
\]
Agora eliminamos o elemento abaixo do segundo pivô
\[
  L_3 \leftarrow L_3 - L_2
\]
Obtemos
\[
\begin{amatrix}{4}
  1 & 2 & -1 & 1  & 3 \\
  0 & 1 & 2  & -1 & 1 \\
  0 & 0 & -1 & 1  & 1
\end{amatrix}
\]
Multiplicando a terceira linha por $-1$
\[
  L_3 \leftarrow - L_3
\]
resulta em
\[
\begin{amatrix}{4}
  1 & 2 & -1 & 1  & 3  \\
  0 & 1 & 2  & -1 & 1  \\
  0 & 0 & 1  & -1 & -1
\end{amatrix}
\]
Eliminamos o elemento acima do terceiro pivô
\[
  L_2 \leftarrow L_2 - 2 L_3
\]
\[
  L_1 \leftarrow L_1 + L_3
\]
Assim
\[
\begin{amatrix}{4}
  1 & 2 & 0 & 0  & 2  \\
  0 & 1 & 0 & 1  & 3  \\
  0 & 0 & 1 & -1 & -1
\end{amatrix}
\]
Agora eliminamos o elemento acima do segundo pivô
\[
  L_1\leftarrow L_1 - 2 L_2
\]
Logo
\[
\begin{amatrix}{4}
  1 & 0 & 0 & -2 & -4 \\
  0 & 1 & 0 & 1  & 3  \\
  0 & 0 & 1 & -1 & -1
\end{amatrix}
\]
A matriz escalonada reduzida representa o sistema
\[
  \systeme[sort={x, y, z, w}]{
    x-2w=-4,
    y+w=3,
    z-w=-1
  }
\]
A variável $w$ é livre pois não possui pivô
\[
  w=t
  \qquad t \in \R
\]
Então
\[
  x=2t-4
\]
\[
  y=3-t
\]
\[
  z=t-1
\]
Portanto, a solução geral do sistema é
\[
  \begin{pmatrix}
    x \\
    y \\
    z \\
    w
  \end{pmatrix}
  =
  \begin{pmatrix}
    -4 \\
    3  \\
    -1 \\
    0
  \end{pmatrix}
  +
  t
  \begin{pmatrix}
    2  \\
    -1 \\
    1  \\
    1
  \end{pmatrix}
  \qquad t \in \R
\]
\end{answer}

% 03
%------------------------------------------------------------------------------%
\exercise{25\%}
Utilize o método de Gauss--Jordan para inverter a matriz
\[
  A =
  \begin{bmatrix}
    2 & 1 & 0 \\
    1 & 1 & 1 \\
    0 & 1 & 1
  \end{bmatrix}
\]

\begin{answer}
Matriz aumentada
\[
  \begin{amatrix}[3]{3}
    2 & 1 & 0 & 1 & 0 & 0 \\
    1 & 1 & 1 & 0 & 1 & 0 \\
    0 & 1 & 1 & 0 & 0 & 1
  \end{amatrix}
\]
Trocamos as linhas $L_1$ e $L_2$
\[
  L_1 \leftrightarrow L_2
\]
Obtemos
\[
\begin{amatrix}[3]{3}
  1 & 1 & 1 & 0 & 1 & 0 \\
  2 & 1 & 0 & 1 & 0 & 0 \\
  0 & 1 & 1 & 0 & 0 & 1
\end{amatrix}
\]
Aplicamos
\[
  L_2 \leftarrow L_2 - 2 L_1
\]
Então
\[
\begin{amatrix}[3]{3}
  1 & 1  & 1  & 0 & 1  & 0 \\
  0 & -1 & -2 & 1 & -2 & 0 \\
  0 & 1  & 1  & 0 & 0  & 1
\end{amatrix}
\]
Fazemos
\[
  L_2 \leftarrow - L_2
\]
Logo
\[
\begin{amatrix}[3]{3}
  1 & 1 & 1 & 0  & 1 & 0 \\
  0 & 1 & 2 & -1 & 2 & 0 \\
  0 & 1 & 1 & 0  & 0 & 1
\end{amatrix}
\]
Agora
\[
  L_3 \leftarrow L_3 - L_2
\]
Obtemos
\[
\begin{amatrix}[3]{3}
  1 & 1 & 1  & 0  & 1  & 0 \\
  0 & 1 & 2  & -1 & 2  & 0 \\
  0 & 0 & -1 & 1  & -2 & 1
\end{amatrix}
\]
Multiplicamos a terceira linha por $-1$
\[
  L_3 \leftarrow - L_3
\]
Assim
\[
\begin{amatrix}[3]{3}
  1 & 1 & 1 & 0  & 1 & 0  \\
  0 & 1 & 2 & -1 & 2 & 0  \\
  0 & 0 & 1 & -1 & 2 & -1
\end{amatrix}
\]
Eliminamos os elementos acima do terceiro pivô
\[
  L_2 \leftarrow L_2 - 2 L_3
\]
\[
  L_1 \leftarrow L_1 - L_3
\]
Resulta
\[
\begin{amatrix}[3]{3}
  1 & 1 & 0 &  1 & -1 &  1 \\
  0 & 1 & 0 &  1 & -2 &  2 \\
  0 & 0 & 1 & -1 &  2 & -1
\end{amatrix}
\]
Finalmente
\[
  L_1 \leftarrow L_1 - L_2
\]
Obtemos
\[
\begin{amatrix}[3]{3}
  1&0&0&0&1&-1\\
  0&1&0&1&-2&2\\
  0&0&1&-1&2&-1
\end{amatrix}
\]
Portanto
\[
  A^{-1} =
  \begin{bmatrix}
    0  & 1  & -1 \\
    1  & -2 & 2  \\
    -1 & 2  & -1
  \end{bmatrix}
\]
\end{answer}

% 04
%------------------------------------------------------------------------------%
\exercise{25\%}
Utilize a inversa para resolver o sistema linear
\[
  \systeme{
    2x + y     = 8,
     x + y + z = 9,
         y + z = 6
  }
\]
\hint{Veja o exercício anterior}

\begin{answer}
A matriz do sistema é
\[
  A =
  \begin{bmatrix}
    2 & 1 & 0 \\
    1 & 1 & 1 \\
    0 & 1 & 1
  \end{bmatrix}
\]
No exercício anterior calculamos
\[
  A^{-1} =
  \begin{bmatrix}
    0  & 1  & -1 \\
    1  & -2 & 2  \\
    -1 & 2  & -1
  \end{bmatrix}
\]
A solução do sistema é
\begin{align*}
  \mathbf{x}
  & = A^{-1} \mathbf{b}
  \\[2mm]
  & =
    \begin{bmatrix}
      0  & 1  & -1 \\
      1  & -2 & 2  \\
      -1 & 2  & -1
    \end{bmatrix}
    \begin{bmatrix}
      8 \\
      9 \\
      6
    \end{bmatrix}
  \\[2mm]
  & =
    \begin{bmatrix}
       0 \cdot 8 + 1 \cdot 9 - 1 \cdot 6 \\
       1 \cdot 8 - 2 \cdot 9 + 2 \cdot 6 \\
      -1 \cdot 8 + 2 \cdot 9 - 1 \cdot 6
    \end{bmatrix}
  \\[2mm]
  & =
    \begin{bmatrix}
      0 +  9 - 6  \\
      8 - 18 + 12 \\
     -8 + 18 - 6
    \end{bmatrix}
  \\[2mm]
  & =
    \begin{bmatrix}
      3 \\
      2 \\
      4
    \end{bmatrix}
\end{align*}
\end{answer}

%------------------------------------------------------------------------------%
\end{document}
%------------------------------------------------------------------------------%

